%% AAAI 2027 Reproducibility Checklist — filled for submission.

\newpage
\section*{Reproducibility Checklist}

\subsubsection*{1. General Paper Structure}

\noindent\textbf{1.1.} Includes a conceptual outline and/or pseudocode description of AI
methods introduced \textit{(yes/partial/no/NA)}: \textbf{partial}

\smallskip
\noindent\textbf{1.2.} Clearly delineates statements that are opinions, hypothesis, and
speculation from objective facts and results \textit{(yes/no)}: \textbf{yes}

\smallskip
\noindent\textbf{1.3.} Provides well-marked pedagogical references for less-familiar readers
to gain background necessary to replicate the paper \textit{(yes/no)}: \textbf{yes}

\subsubsection*{2. Theoretical Contributions}

\noindent\textbf{2.1.} Does this paper make theoretical contributions? \textit{(yes/no)}:
\textbf{no}

\smallskip
\noindent\textit{Items 2.2--2.8 are not applicable (no theoretical contributions).}

\subsubsection*{3. Dataset Usage}

\noindent\textbf{3.1.} Does this paper rely on one or more datasets? \textit{(yes/no)}: \textbf{yes}

\smallskip
\noindent\textbf{3.2.} A motivation is given for why the experiments are conducted on the
selected datasets \textit{(yes/partial/no/NA)}: \textbf{yes}

\smallskip
\noindent\textbf{3.3.} All novel datasets introduced in this paper are included in a data
appendix \textit{(yes/partial/no/NA)}: \textbf{NA}

\smallskip
\noindent\textbf{3.4.} All novel datasets introduced in this paper will be made publicly
available upon publication with a license that allows free usage for research purposes
\textit{(yes/partial/no/NA)}: \textbf{NA}

\smallskip
\noindent\textbf{3.5.} All datasets drawn from the existing literature are accompanied by
appropriate citations \textit{(yes/no/NA)}: \textbf{yes}

\smallskip
\noindent\textbf{3.6.} All datasets drawn from the existing literature are publicly available
\textit{(yes/partial/no/NA)}: \textbf{partial}

\smallskip
\noindent\textbf{3.7.} All datasets that are not publicly available are described in detail,
with explanation why publicly available alternatives are not scientifically satisficing
\textit{(yes/partial/no/NA)}: \textbf{partial}

\subsubsection*{4. Computational Experiments}

\noindent\textbf{4.1.} Does this paper include computational experiments? \textit{(yes/no)}:
\textbf{yes}

\smallskip
\noindent\textbf{4.2.} This paper states the number and range of values tried per
(hyper-)parameter during development, along with the criterion used for selecting the final
parameter setting \textit{(yes/partial/no/NA)}: \textbf{partial}

\smallskip
\noindent\textbf{4.3.} Any code required for pre-processing data is included in the appendix
\textit{(yes/partial/no)}: \textbf{partial}

\smallskip
\noindent\textbf{4.4.} All source code required for conducting and analyzing the experiments
is included in a code appendix \textit{(yes/partial/no)}: \textbf{partial}

\smallskip
\noindent\textbf{4.5.} All source code required for conducting and analyzing the experiments
will be made publicly available upon publication with a license that allows free usage for
research purposes \textit{(yes/partial/no)}: \textbf{yes}

\smallskip
\noindent\textbf{4.6.} All source code implementing new methods have comments detailing the
implementation, with references to the paper where each step comes from
\textit{(yes/partial/no)}: \textbf{partial}

\smallskip
\noindent\textbf{4.7.} If an algorithm depends on randomness, the method used for setting seeds
is described in a way sufficient to allow replication of results
\textit{(yes/partial/no/NA)}: \textbf{yes}

\smallskip
\noindent\textbf{4.8.} This paper specifies the computing infrastructure used for running
experiments (hardware and software), including GPU/CPU models; amount of memory; OS; names and
versions of relevant software libraries and frameworks \textit{(yes/partial/no)}: \textbf{yes}

\smallskip
\noindent\textbf{4.9.} This paper formally describes evaluation metrics used and explains the
motivation for choosing these metrics \textit{(yes/partial/no)}: \textbf{yes}

\smallskip
\noindent\textbf{4.10.} This paper states the number of algorithm runs used to compute each
reported result \textit{(yes/no)}: \textbf{yes}

\smallskip
\noindent\textbf{4.11.} Analysis of experiments goes beyond single-dimensional summaries of
performance to include measures of variation, confidence, or other distributional information
\textit{(yes/no)}: \textbf{yes}

\smallskip
\noindent\textbf{4.12.} The significance of any improvement or decrease in performance is
judged using appropriate statistical tests (e.g., Wilcoxon signed-rank)
\textit{(yes/partial/no)}: \textbf{yes}

\smallskip
\noindent\textbf{4.13.} This paper lists all final (hyper-)parameters used for each
model/algorithm in the paper's experiments \textit{(yes/partial/no/NA)}: \textbf{yes}
